\section{PWM DC/DC converters in CCM}

\subsection{Ctrl+F keywords: CCM, continuous conduction, buck, boost, buck-boost, Cuk, volt-second, small-ripple, duty ratio, ripple voltage, component rating, SMPS}

Ideal switches, no losses, large \(C\) so \(V_o\) is DC plus small ripple. Small-ripple: replace instantaneous \(v_o\) by \(V_o\) in inductor volt-second equations.

\subsubsection{Stored energy, switching frequency and filter cutoff}
\[
E_L=\frac12 L I^2,\qquad E_C=\frac12 C V^2,\qquad P=\frac{E}{T_s}=E f_s.
\]
Choose LC-filter cutoff \(f_c\ll f_s\) so switching harmonics are rejected:
\[
f_c=\frac{1}{2\pi\sqrt{LC}}.
\]

\subsubsection{Duty ratio from PWM}
\[
D=\frac{t_{\mathrm{on}}}{T_s},\qquad 0<D<1.
\]

\subsubsection{Volt-second balance (CCM)}
Over one period the net inductor volt-seconds is zero:
\[
\int_0^{T_s} v_L\,\mathrm{d}t=0.
\]
This is the fastest route to \(M=V_o/V_{\mathrm{in}}\) in CCM. Charge balance on capacitors gives current ratios.

\subsection{CCM conversion ratios}
\begin{center}
\begin{tabular}{@{}l l l l@{}}
\toprule
Topology & \(M=V_o/V_{\mathrm{in}}\) & Range & Notes\\
\midrule
Buck & \(D\) & \(0\le M\le 1\) & Step-down, continuous input chopped\\
Boost & \(1/(1-D)\) & \(M\ge 1\) & \(M\to\infty\) as \(D\to 1\) only if lossless\\
Buck-boost & \(D/(1-D)\) & \(0\le M<\infty\) & Inverted \(V_o\); product of buck and boost\\
Cuk & \(D/(1-D)\) & \(0\le M<\infty\) & Inverted; continuous \(i_{L1}\) and \(i_{L2}\)\\
\bottomrule
\end{tabular}
\end{center}
Losses cap boost/Cuk/buck-boost gain as \(D\to 1\). \(D=1\) is not a usable Cuk operating point (\(V_o=0\) because the transistor never releases energy).

\subsubsection{Buck CCM: inductor ripple from volt-seconds}
\[
\Delta i_L=\frac{V_o(1-D)T_s}{L}=\frac{(V_{\mathrm{in}}-V_o)DT_s}{L},\qquad V_o=D V_{\mathrm{in}}.
\]
On/off inductor currents (average \(\overline{i}_L=I_o\)):
\[
i_{L,\mathrm{on}}(t)=-\frac{\Delta i_L}{2}+\frac{\Delta i_L}{D T_s}t+\overline{i}_L,\qquad
i_{L,\mathrm{off}}(t)=\frac{\Delta i_L}{2}-\frac{\Delta i_L}{(1-D)T_s}t+\overline{i}_L.
\]

\subsubsection{Buck CCM: capacitor current and output ripple}
With \(i_C=i_L-I_o\),
\[
i_{C,\mathrm{rms}}=\frac{\Delta i_L}{\sqrt{12}}=\frac{I_o(1-D)}{\sqrt{3K}},\qquad K=\frac{2L}{R T_s}.
\]
Peak-to-peak output ripple (triangle area on \(C\)):
\[
\Delta V_o=\frac{V_o(1-D)T_s^2}{8LC}=\frac{V_o(1-D)\pi^2}{2}\left(\frac{f_c}{f_s}\right)^2.
\]

\subsubsection{Boost CCM: conversion and output ripple}
\[
V_{\mathrm{in}} D T_s=(V_o-V_{\mathrm{in}})(1-D)T_s \quad\Rightarrow\quad \frac{V_o}{V_{\mathrm{in}}}=\frac{1}{1-D}.
\]
Capacitor supplies \(I_o\) while the diode is off:
\[
\Delta V_o=\frac{I_o D T_s}{C}=\frac{D V_o T_s}{RC},\qquad \frac{\Delta V_o}{V_o}=\frac{D T_s}{RC}.
\]
Boost (and buck-boost) output ripple is load-dependent and independent of \(L\). Capacitor current stress is harsher than in a buck.

\subsubsection{Buck-boost CCM: currents and conversion}
\[
\Delta i_L=\frac{V_{\mathrm{in}} D T_s}{L}=\frac{V_o(1-D)T_s}{L},\qquad \frac{V_o}{V_{\mathrm{in}}}=\frac{D}{1-D}.
\]
Average diode current is \(I_o\); average transistor current is \(I_{\mathrm{in}}\):
\[
I_o=\overline{i}_L(1-D),\qquad I_{\mathrm{in}}=\overline{i}_L D,\qquad \overline{i}_L=I_{\mathrm{in}}+I_o=\frac{I_o}{1-D}.
\]
Output ripple (same construction as boost):
\[
\frac{\Delta V_o}{V_o}=\frac{D T_s}{RC}.
\]

\subsubsection{Cuk CCM: transfer capacitor and conversion}
Volt-seconds on \(L_1\) and \(L_2\):
\[
V_{C1}=\frac{V_{\mathrm{in}}}{1-D}=\frac{V_o}{D}=V_{\mathrm{in}}+V_o,\qquad \frac{V_o}{V_{\mathrm{in}}}=\frac{D}{1-D}.
\]
Current ratio from power balance:
\[
\frac{I_o}{I_{\mathrm{in}}}=\frac{\overline{i}_{L2}}{\overline{i}_{L1}}=\frac{1-D}{D}.
\]
Ripples:
\[
\Delta i_{L1}=\frac{V_{\mathrm{in}} D T_s}{L_1}=\frac{V_o(1-D)T_s}{L_1},\qquad
\Delta i_{L2}=\frac{V_{\mathrm{in}} D T_s}{L_2}=\frac{V_o(1-D)T_s}{L_2},
\]
\[
\Delta V_{C2}=\frac{V_o(1-D)T_s^2}{8 C_2 L_2}.
\]

\subsection{CCM rms of a DC current plus triangular ripple}
\[
I_{\mathrm{rms}}^2=\overline{I}^2+\frac{(\Delta i)^2}{12}.
\]
A switch that carries that waveform only for fraction \(D\) has
\[
I_{\mathrm{rms}}=\sqrt{D}\sqrt{\overline{I}^2+\frac{(\Delta i)^2}{12}}.
\]

\subsection{CCM component ratings (ideal)}
Blocking voltages: buck \(V_{\mathrm{in}}\); boost \(V_o\); buck-boost and Cuk \(V_{\mathrm{in}}+V_o\). Average transistor current \(I_{\mathrm{in}}\) (Cuk: \(I_{\mathrm{in}}\); switch sees \(I_{\mathrm{in}}+I_o\) during on). Average diode current \(I_o\). Capacitor average current is zero in periodic steady state.

\begin{center}
\begin{tabular}{@{}llll@{}}
\toprule
 & Buck & Boost & Buck-boost / Cuk switch-diode\\
\midrule
\(T\) voltage & \(V_{\mathrm{in}}\) & \(V_o\) & \(V_{\mathrm{in}}+V_o\)\\
\(D\) voltage & \(V_{\mathrm{in}}\) & \(V_o\) & \(V_{\mathrm{in}}+V_o\)\\
\(T\) average & \(I_o D=I_{\mathrm{in}}\) & \(I_{\mathrm{in}} D\) & \(I_{\mathrm{in}}\)\\
Diode average & \(I_o(1-D)\) & \(I_o\) & \(I_o\)\\
\bottomrule
\end{tabular}
\end{center}
Cuk extras: \(C_1\) sees \(V_{\mathrm{in}}+V_o\); \(C_2\) sees \(V_o\); \(L_1\) average \(I_{\mathrm{in}}\); \(L_2\) average \(I_o\). Inductor voltage rating in exercises is taken as \(\max(V_{\mathrm{in}},V_o)\).

\subsection{Exam recipe: buck CCM output ripple}
\textbf{Use when:} given \(V_{\mathrm{in}}\), \(D\) or \(V_o\), \(L\), \(C\), \(f_s\); asked for \(\Delta V_o\).

\textbf{Steps:}
\begin{enumerate}
  \item \(V_o=D V_{\mathrm{in}}\) if \(V_o\) is not given.
  \item \(f_c=1/(2\pi\sqrt{LC})\).
  \item \(\Delta V_o=V_o\dfrac{(1-D)\pi^2}{2}(f_c/f_s)^2\), or equivalently \(\Delta V_o=V_o(1-D)T_s^2/(8LC)\).
\end{enumerate}

\textbf{Fast checks:} \(\Delta V_o\) falls if \(D\) rises (fixed \(V_{\mathrm{in}}\), \(L\), \(C\), \(f_s\)); \(\Delta V_o=0\) at \(D=1\).

\textbf{Common traps:} using boost's \(D T_s/(RC)\) ripple on a buck; forgetting \(V_o=DV_{\mathrm{in}}\) before inserting into \(\Delta V_o/V_o\).

\subsection{Exam recipe: Cuk CCM voltage ratings}
\textbf{Use when:} Cuk, CCM, given \(f_s\), \(t_{\mathrm{on}}\), \(V_{\mathrm{in}}\); asked for device/capacitor voltage ratings.

\textbf{Steps:}
\begin{enumerate}
  \item \(D=t_{\mathrm{on}} f_s=t_{\mathrm{on}}/T_s\).
  \item \(V_o=D V_{\mathrm{in}}/(1-D)\) (magnitude; polarity inverted).
  \item Transistor, diode and \(C_1\): \(V_{\mathrm{in}}+V_o\). \(C_2\), \(L_1\), \(L_2\): \(V_o\) (course table uses \(V_o\) for the two inductors).
\end{enumerate}

\textbf{Fast checks:} \(V_{C1}=V_{\mathrm{in}}/(1-D)=V_o/D\).

\textbf{Common traps:} rating the switch at \(V_{\mathrm{in}}\) only; using DCM formulae while the problem states CCM.

\subsection{Recognition rules}
CCM if the inductor (or both Cuk inductors) never returns to zero in the period. Volt-second \(\Rightarrow\) load-independent \(M(D)\). Input current continuous on boost and Cuk \(L_1\); discontinuous on buck (chopped) and on both ports of buck-boost.

\subsection{Common traps}
Signs: buck-boost and Cuk invert. Do not set \(D=1\) for infinite boost gain. Filter \(f_c\) is not \(f_s\). \(K=2L/(RT_s)\) is not needed for CCM gain, only for boundary and some rms forms.

\subsection{One-line checks}
Buck: \(V_o<V_{\mathrm{in}}\). Boost: \(V_o>V_{\mathrm{in}}\). Buck-boost/Cuk: \(M=1\) at \(D=0.5\). Power: \(V_{\mathrm{in}}I_{\mathrm{in}}=V_o I_o\) if lossless.
